\section{Unique Representation Theorem} \label{sec-unique-repr-thm}

In this chapter, we present the main theorem and complete the proofs given by \textcite[pp. 195-199]{foundations2}. The remainder of this chapter, except for the final section, consists of the proof of this theorem. The last section is dedicated to proving an intermediate result from functional analysis that is used in the proof.\\
We begin by stating the theorem, and follow a top-down approach in the proof to provide a broad understanding. This means, we use results established in \cref{sect-homogeneity-theorem}. The same approach applies to the proofs in that section; see \cref{sect-homogeneity-theorem} for further explanation.
We start with the statement of the theorem.

\begin{thm}[Unique Representation Theorem: The Power Metric as the Sole Representation for Segmentally Additive AD Structures]\label{thm-uniquene-repr-minkowski} \index{Theorem!unique representation}
	Suppose $\langle \ps{A}, \succsim \rangle$ is a complete, segmentally additive, factorial proximity structure, that is also and additive difference structure.
	Then there exists a unique $p \geq 1$ and real-valued functions $\varphi_i$, defined on $\ps{A}_i$, $i = 1, \ldots, n$, that are unique up to a positive linear transformation, such that, for all $\ps{a}, \ps{b}, \ps{c}, \ps{d} \in \ps{A}$,
	$$(\ps{a}, \ps{b}) \succsim (\ps{c}, \ps{d}) \quad  \text{ iff } \quad  \delta(\ps{a}, \ps{b}) \geq \delta(\ps{c}, \ps{d}),$$
	where
	$$\delta(\ps{a}, \ps{b}) = \left[\sum_{i=1}^{n} |\varphi_i(\ps{a}) - \varphi_i(\ps{b})|^p\right]^{1/p}.$$
\end{thm}

\newcommand\ta{x}

\begin{proof}
	This proof proceeds in multiple steps.
	First, we make observations about the domains and ranges of the functions $F_i$, use this to find bounds for the domain of $F$. Using these bounds, we argue that $F$ can be extended to the positive reals.
	Next, we derive from this a functional equation, that is only satisfied by $f(x) = \alpha x^p$.
	Finally, we prove that $F$ is equal to this extension on the whole domain, and therefore $\metric$ is a Minkowski metric.
	\setcounter{proofpart}{-1}
	\begin{proofpart}[Establishing the basics and using the \cref{thm-homogeneity}]
		By \Cref{thm-conds} and \cref{rem-metric-form}, $\metric$ must be of the form
		\begin{equation*}
			\metr{\ps{a}}{\ps{b}} = F^{-1} \left [ \sum_{i=1}^{n} F(| \varphi_i(\ps{a}_i) - \varphi_i(\ps{b}_i)|)  \right]
		\end{equation*}
		where $F$ is continuous, convex, and increasing.
		By the same argument as in \crefpart{lemma-rem}{rem-lemma-domains}, $\metric$ must be defined on a Cartesian product of intervals of the form $[0,\rho_i)$ where $\rho_i$ are the ranges of $| \varphi_i(\ps{a}_i) - \varphi_i(\ps{b}_i)|$ for $i=1, \dots, n$. This Cartesian product is convex and, by leaving out $0$ in each product, also open. Therefore, the requirements for \cref{thm-homogeneity} are met if we leave out $\ve{0}$ and applying the lemma yields homogeneity for $\ve{x}, \ve{z} \neq \ve{0}$:
    %METR
		$$\metr{\ve{x}}{\ve{x}+ t(\ve{z}-\ve{x})} = t \metr{\ve{x}}{\ve{z}}.$$
		Applying homogeneity for the definition of $\metric$ above and setting $\ta_i = |\varphi_i(\ps{b}_i) - \varphi_i(\ps{a}_i)|$ yields
		\begin{equation}\label{proof-homogen-eq}
			F^{-1} \left [ \sum_{i=1}^{n} F(t \ta_i)  \right] = t \cdot F^{-1} \left [ \sum_{i=1}^{n} F(\ta_i)  \right].
		\end{equation}
		However, due to translation invariance, homogeneity also holds for $\ve{x} = \ve{0}$\footnote{By adding $0 = \ve{y} - \ve{y}$ for some $\ve{y}$ in the domain, we get homogeneity for $\ve{0}$:
			\begin{equation*}
				\metr{\ve{0}}{t \ve{z}} =F^{-1} \left [ \sum_{i=1}^{n} F(t (y_i - (y_i - z_i))) \right]
				= t F^{-1} \left [ \sum_{i=1}^{n} F(y_i - (y_i - z_i)) \right]
				= t \metr{\ve{0}}{\ve{z}}.
			\end{equation*}}. We will show, that the only solutions to this equation are of the form $F(x) = \alpha x^p$, i.e., the power functions.

		First, we start by observing the domains and the ranges of $F$.
		By \cref{thm-conds}, the additive difference model
		\begin{equation*}
			\delta(\ps{a}, \ps{b}) = G \left( \sum_{i=1}^{n} F_i \left( |\varphi_i(\ps{a}_i) - \varphi_i(\ps{b}_i)| \right) \right)
		\end{equation*}
		yields a metric (that is additive on the coordinate axes) if $F_i(\ta) = F(t_i \ta)$ and $F = G^{-1}$ for $i=1, ..., n$. Thus, $F^{-1}$ is defined for all numbers $\sum_{i=1}^{n} F_i \left( |\varphi_i(\ps{a}_i) - \varphi_i(\ps{b}_i)| \right)$, and each $F_i$ coincides with $F$ on its domain (except for the factor $t_i$, which scales the argument and adjusts the domain).
		Since the domains of $F_i$ are intervals (because $\delta$ is a metric with additive segments and therefore for any point in the domain of $F_i$, the interval between zero and that point must also be in the domain), we can choose $\omega > 0$ such that for each $i = 1, \ldots, n$ and any $\ta$ with $0 \leq \ta \leq \omega/t_i$, $\ta$ lies within the domain of $F_i$.
		In particular, if we set $\omega_i = \omega/t_i$, then $F^{-1} \left[ \sum_{i=1}^{n} F_i(\omega_i) \right] = F^{-1} \left[ nF(\omega) \right] =: \Omega$ is defined. Thus, $F = G^{-1}$ is defined at least on the interval $[0, \Omega]$. We can now extend $F$ as follows.
	\end{proofpart}
	\begin{proofpart}[Extending the bounds of $F$]
		For $\ta < \omega$, we have by \Cref{proof-homogen-eq},
		\begin{align*}
			F^{-1}[nF(\ta)] & = F^{-1}[nF(\ta \omega/\omega)]   \\
			                & = (\ta/\omega) F^{-1}[nF(\omega)] \\
			                & = \ta \Omega/\omega.
		\end{align*}
		Thus, $F$ satisfies the following, by plugging in the correct values for $\ta$ ($\ta \frac{\omega}{\Omega}$ and $F^{-1}( \frac{\ta}{n} )$)
		\begin{align}
			F(\ta)      & = nF(\ta \omega/\Omega), \quad 0 \leq \ta \leq \Omega; \label{proof-f-erw}               \\
			F^{-1}(\ta) & = (\Omega/\omega) F^{-1}(\ta/n), \quad 0 \leq \ta \leq nF(\omega). \label{proof-f-1-erw}
		\end{align}
		These two equations can now be used to extend the domain of $F$ and $F^{-1}$.

		We use \Cref{proof-f-erw} by defining $F$ through the right-hand side (for values where the right-hand side is defined, i.e., for $0 \leq \ta \leq \Omega^2/\omega$). Since $\Omega^2/\omega > \Omega$, this indeed extends $F$ beyond the interval $[0, \Omega]$.
		Similarly, $F^{-1}$ satisfies \Cref{proof-f-1-erw} for the extended interval $0 \leq \ta \leq nF(\Omega) = n^2 F(\omega)$.
		Moreover, \Cref{proof-homogen-eq} is now valid for $0 \leq t \leq 1$ and $0 \leq \ta_i \leq \Omega$:
		{\allowdisplaybreaks[1]
      \begin{align*}
			F^{-1} \left[ \sum_{i=1}^{n} F(t \ta_i) \right] & = F^{-1} \left[ \sum_{i=1}^{n} nF(t \ta_i \omega / \Omega) \right] \tag{by \cref{proof-f-erw}}                                                               \\
			                                                & = F^{-1} \left[ n \sum_{i=1}^{n} F(t \ta_i \omega / \Omega) \right]                                                                                          \\
			                                                & = \frac{\Omega}{\omega} F^{-1} \left[ \sum_{i=1}^{n} F(t \ta_i \omega / \Omega) \right] \tag{by \cref{proof-f-1-erw}}                                        \\
			                                                & = \frac{\Omega}{\omega} t F^{-1} \left[ \sum_{i=1}^{n} F(\ta_i \omega / \Omega) \right] \tag{by homogeneity applied for $\ta_i \omega / \Omega \leq \omega$} \\
			                                                & = t F^{-1} \left[ n \sum_{i=1}^{n} F(\ta_i \omega / \Omega) \right] \tag{by \cref{proof-f-1-erw}}                                                            \\
			                                                & = t F^{-1} \left[ \sum_{i=1}^{n} n F(\ta_i \omega / \Omega) \right]                                                                                          \\
			                                                & = t F^{-1} \left[ \sum_{i=1}^{n} F(\ta_i) \right]. \tag{by \cref{proof-f-erw}}
		\end{align*}}
		Further, the extended function is both increasing and continuous since the extension is based on scaling of the domain and the range of the original function.

		For continuity let $\ta_k \rightarrow \ta$ with $\omega < \ta < \Omega$. Then
		\begin{align*}
			\lim_{k \rightarrow \infty} F(\ta_k) & = \lim_{k \rightarrow \infty} n F \left ( \ta_k \frac{\omega}{\Omega} \right) \tag{Definition of $F(t)$ for $\omega \leq t \leq \Omega$} \\
			                                     & = n F \left ( \lim_{k \rightarrow \infty} \ta_k \frac{\omega}{\Omega} \right) \tag{Continuity of $F(t)$ for $0 \leq t \leq \omega$}      \\
			                                     & = n F \left ( \ta \frac{\omega}{\Omega} \right) \tag{Convergence of $\ta_k$}                                                             \\
			                                     & = F(\ta) \tag{Definition of $F(t)$ for $\omega \leq t \leq \Omega$}
		\end{align*}
		which is equivalent to $F$ begin continuous on the extended domain.

		For monotonicity, we can use a similar argument. Let $\ta < \ta'$ with $\omega < \ta, \ta' < \Omega$. Then $\ta \frac{\omega}{\Omega} < \ta' \frac{\omega}{\Omega}$ holds. Plugging that into $F$ and multiplying by $n$ yields $nF \left(\ta \frac{\omega}{\Omega}\right) <nF \left( \ta' \frac{\omega}{\Omega}\right)$ due to $F$ being increasing on the original domain. Therefore, $\ta < \ta'$ implies $F(\ta) < F(\ta')$ which is equivalent to $F$ increasing.
		Thus, we have established strict increasing monotonicity, continuity and \Cref{proof-homogen-eq} for the extended function.

		Repeatedly applying the arguments, extends the intervals in which $F$ is defined and in which monotonicity, continuity, and homogeneity (\cref{proof-homogen-eq}) are valid by a factor of $\Omega/\omega$ each time. Thus, we can extend $F$ to $[0, \infty)$ while preserving these properties.

		We can prove further, that the homogeneity equation must also hold for $t \in [0, \infty)$.	For that, let $\ta_i > 0$ and $t > 1$. Then there exists $t' = \frac{1}{t} < 1$ such that $t' \cdot t \ta_i = \ta_i$. Therefore,
		\begin{align*}
			t t' F^{-1} \left[ \sum_{i=1}^{n} F(\ta_i) \right] & = F^{-1} \left[ \sum_{i=1}^{n} F(\ta_i) \right] \tag{$t t' = 1$}     \\
			                                                   & = F^{-1} \left[ \sum_{i=1}^{n} F(t t'\ta_i) \right] \tag{$t t' = 1$} \\
			                                                   & = t' F^{-1} \left[ \sum_{i=1}^{n} F(t \ta_i) \right]
		\end{align*}
		The last step holds, because both $t'<1$ and $t \ta_i < \infty$ hold.
		Dividing both sides by $t'$ yields
		\begin{equation*}
			F^{-1} \left[ \sum_{i=1}^{n} F(t \ta_i) \right] = t F^{-1} \left[ \sum_{i=1}^{n} F(\ta_i) \right] \quad \text{ for } \ta, t < \infty .
		\end{equation*}
		To sum up, we have obtained an extended function which satisfies monotonicity, continuity, and homogeneity for $t \in [0, \infty)$. This extended function must coincide with tie original $F$ on at least the interval $[0, \Omega]$ (because this was the domain of $F$ that we deduced, before extending it).
	\end{proofpart}

	\begin{proofpart}[Deriving functional equation]
		Next, we derive a functional equation that allows us to use results on quasi-arithmetic homogeneous means. The theorem states as follows:
		\begin{thm*}[Uniqueness of homogeneous quasi-arithmetic means, \cref{thm-power-mean-uniqueness}]
			Suppose that $\phi$ is continuous and strictly increasing in the open interval $(0,\infty)$ and that
			\begin{equation*}
				\mean{t \ve{x}} = t \mean{ \ve{x}} \quad  \text{with} \quad \mean{\ve{x} } = \phi^{-1} \left ( \sum_{i=1}^{n} w_i \phi(x_i) \right)
			\end{equation*}
			for all positive $\ve{x}, \ve{w}, t$ and all $n \in \mathbb{N}$ such that $\sum_{i=1}^{n} w_i = 1$. Then $\mean{\ve{x}} = \meanp{\ve{x}}$. (i.e., $\mean{\ve{x}}$ is either a power mean with $\phi(x) = \alpha x^p + \beta$ and $\alpha, \beta, p \in \mathbb{R}, p > 0$ or a geometric mean with $\phi(x) = \alpha \log (x)$, $\alpha \in \mathbb{R}$.)
		\end{thm*}
		The proof of this statement will be presented in \cref{sect-HLP_proof}, here we will apply the theorem for our proof.

		Since all requirements for $\phi$ are already met by $F$, we have to derive the functional equation for all $n \in \mathbb{N}$ and all positive weights. First, we will prove that the homogeneity holds not only for sums of $n$ summands, but for arbitrary many summands $N \in \mathbb{N}$. Then we will incorporate equal weights into this equation, extend them to commensurable weights (i.e., rational weights with a common denominator) and then use this to approximate arbitrary real weights.

		\begin{proofpp}[Homogeneity for arbitrary long sums]
			Let for $N < n$ the first $N$ coordinates of $\ve{x}$ be zero, i.e., $\ve{x} = (\ta_1, \dots, \ta_N, 0, \dots , 0)$. Then, since $F(0) = 0$ it holds that, $F(x_i) = 0$ for $i > N$. Applying homogeneity for this $\ve{x}$ yields
			\begin{align}
				F^{-1} \left[ \sum_{i=1}^{N} F(t \ta_i) \right] & = F^{-1} \left[ \sum_{i=1}^{n} F(t \ta_i) \right] \tag{$x_i = 0$ for $i > N$}  \\
				                                                & = t F^{-1} \left[ \sum_{i=1}^{n} F(\ta_i) \right] \tag{homogeneity}            \\
				                                                & = t F^{-1} \left[ \sum_{i=1}^{N} F(\ta_i) \right]. \tag{$x_i = 0$ for $i > N$}
			\end{align}
			Therefore, we have homogeneity for $N \leq n$:
			\begin{equation} \label{proof-homogen-ksmall}
				F^{-1} \left[ \sum_{i=1}^{N} F(t \ta_i) \right] = t F^{-1} \left[ \sum_{i=1}^{N} F(\ta_i) \right].
			\end{equation}
			Now, we proceed with homogeneity for $n+1$ summands. This can be achieved by grouping the last two summands into one and applying homogeneity twice:
			{\allowdisplaybreaks[1] % Applies only to the next align environment
			\begin{align*}
				F^{-1} \left[ \sum_{i=1}^{n+1} F(t\ta_i) \right] & = F^{-1} \left[ \sum_{i=1}^{n-1} F(t\ta_i) + F(t\ta_{n}) + F(t\ta_{n+1}) \right]      \tag{expanding sum}                                                             \\
				                                                 & = F^{-1} \left[ \sum_{i=1}^{n-1} F(t\ta_i) + F (F^{-1} ( F(t\ta_{n}) + F(t\ta_{n+1}) )) \right] \tag{regrouping}                                                      \\
				                                                 & =  F^{-1} \left[ \sum_{i=1}^{n-1} F(t\ta_i) + F (t F^{-1} ( F(\ta_{n}) + F(\ta_{n+1}) )) \right] \tag{\small{\cref{proof-homogen-ksmall} on the last summand, $N=2$}} \\
				                                                 & = t F^{-1} \left[ \sum_{i=1}^{n-1} F(\ta_i) + F ( F^{-1} ( F(\ta_{n}) + F(\ta_{n+1}) )) \right] \tag{\small{homogeneity for $n$ summands}}                            \\
				                                                 & = t F^{-1} \left[ \sum_{i=1}^{n-1} F(\ta_i) + F(\ta_{n}) + F(\ta_{n+1}) \right] \tag{\small{cancellation of inverses}}                                                \\
				                                                 & = t F^{-1} \left[ \sum_{i=1}^{n+1} F(\ta_i) \right].
			\end{align*}
			}
			Therefore, we now have homogeneity for $n+1$ summands. Repeating this, yields homogeneity for arbitrary many summands:
			\begin{equation} \label{proof-homogen-kbig}
				F^{-1} \left[ \sum_{i=1}^{N} F(t \ta_i) \right] = t F^{-1} \left[ \sum_{i=1}^{N} F(\ta_i) \right] \quad N \in \mathbb{N}.
			\end{equation}
		\end{proofpp}

		\begin{proofpp}[Homogeneity with weights]
			Now, we proceed with incorporating weights into the equation. Firstly, applying \Cref{proof-f-1-erw} on the left-hand side of \Cref{proof-homogen-kbig} yields
			\begin{equation*}
				\Omega / \omega F^{-1} \left[ \frac{1}{n} \sum_{i=1}^{N} F(t \ta_i) \right] = F^{-1} \left[ \sum_{i=1}^{N} F(t \ta_i) \right] = t F^{-1} \left[ \sum_{i=1}^{N} F(\ta_i) \right] .
			\end{equation*}
			Applying \Cref{proof-f-1-erw} once more on the right-hand side, yields
      {\small \begin{equation*}
				\Omega / \omega F^{-1} \left[ \frac{1}{n} \sum_{i=1}^{N} F(t \ta_i) \right] = F^{-1} \left[ \sum_{i=1}^{N} F(t \ta_i) \right] = t F^{-1} \left[ \sum_{i=1}^{N} F(\ta_i) \right] =\Omega / \omega t F^{-1} \left[ \frac{1}{n} \sum_{i=1}^{N} F(\ta_i) \right].
    \end{equation*}}
			Hence,
			\begin{equation*}
				F^{-1} \left[ \frac{1}{n} \sum_{i=1}^{N} F(t \ta_i) \right] = t F^{-1} \left[ \frac{1}{n} \sum_{i=1}^{N} F(\ta_i) \right]
			\end{equation*}
			must hold for all $N \in \mathbb{N}$. Repeatedly applying this argument yields that
			\begin{equation} \label{proof-homogen-prower-weights}
				F^{-1} \left[ \frac{1}{n^l} \sum_{i=1}^{N} F(t \ta_i) \right] = t F^{-1} \left[ \frac{1}{n^l} \sum_{i=1}^{N} F(\ta_i) \right]
			\end{equation}
			must hold for all $l, N \in \mathbb{N}$.

			Now we extend this to weights of the form $\frac{q}{n^l}$ with $q, l \in \mathbb{N}$. For this, denote the set of all these numbers as $\mathbb{W} := \{\frac{q}{n^l}: q, l \in \mathbb{N}\}$. Let $\ve{w} \in \mathbb{W}^N$. Then we can write $w_i = \frac{q_i}{n^{l_i}}$ for all $i = 1, \dots, N$. We can assume that $l_i$ is the same for all $i$ without loss of generality, by having $n_l$ be a common denominator and expanding all fractions accordingly. Therefore, we not have weights $w_i = \frac{q_i}{n^l}$ for $i=1, \dots N$.
			The basic idea now is, to expand each summand $q_i \cdot F(t x_i)$ into a sum of $F(t x_i)$ with length $q_i$.
			For that, let $Q = \sum_{i=1}^{N} q_i$ and define $\tilde \ta_i$ by repeating $x_i$ for $q_i$ times:
			$$\tilde \ta_i =
				\begin{cases}
					x_1 & i=1, \dots, q_1         \\
					x_2 & i=q_1+1, \dots, q_2     \\
					    & \vdots                  \\
					x_N & i=q_{N-1}+1, \dots, q_N
				\end{cases}.$$
			Then
			\begin{align*}
				F^{-1} \left[ \sum_{i=1}^{N} w_i F(t \ta_i) \right] & = F^{-1} \left[ \sum_{i=1}^{N} \frac{q_i}{n^l} F(t \ta_i) \right] \tag{plugging in $w_i = \frac{q_i}{n^l}$ }     \\
				                                                    & = F^{-1} \left[ \frac{1}{n^l} \sum_{i=1}^{N} q_i F(t \ta_i) \right] \tag{rearranging}                            \\
				                                                    & = F^{-1} \left[ \frac{1}{n^l} \sum_{i=1}^{N} \sum_{j=1}^{q_i} F(t \ta_i) \right] \tag{$q_i \in \mathbb{N}$}      \\
				                                                    & = F^{-1} \left[ \frac{1}{n^l} \sum_{i=1}^{Q} F(t \tilde \ta_i) \right] \tag{Definition of $\tilde \ta_i$}        \\
				                                                    & = t F^{-1} \left[ \frac{1}{n^l} \sum_{i=1}^{Q} F(\tilde \ta_i) \right] \tag{\cref{proof-homogen-prower-weights}} \\
				                                                    & = t F^{-1} \left[ \sum_{i=1}^{N} w_i F(\ta_i) \right] \tag{\small{reversing the steps}}
			\end{align*}
			Thus, we have derived
			\begin{equation} \label{proof-homogen-rats}
				F^{-1} \left[ \sum_{i=1}^{N} w_i F(t \ta_i) \right] = t F^{-1} \left[ \sum_{i=1}^{N} w_i F(\ta_i) \right] \quad w_i \in \mathbb{W}.
			\end{equation}

			Now, in order to approximate real weights, we need density of $\mathbb{W}$ in $\mathbb{R}_{\geq 0}$. Since $n^l$ gets arbitrarily small, there exists for every $\varepsilon >0$ an $l$ such that $n^l < \varepsilon$. Therefore, if $x - y < \varepsilon$, for $x, y \in \mathbb{R}_{\geq 0}$ and $x > y$ there exists $k \in \mathbb{N}$ such that $x < \frac{k}{n^l} < y$. Thus, $\mathbb{W}$ is dense in $\mathbb{R}_{\geq 0}$, and we can approximate real weights with converging sequences from $\mathbb{W}$.

			Finally, let $w_i \in \mathbb{R}_{\geq 0}$ and $q_{i_k} \in \mathbb{W}$ be sequences for $i=1, \dots, N$ such that $q_{i_k} \rightarrow w_i$ for $k \rightarrow \infty$. Then
			\begin{align*}
				F^{-1} \left[ \sum_{i=1}^{N} w_i F(t \ta_i) \right] & = F^{-1} \left[ \sum_{i=1}^{N} \lim_{k \rightarrow \infty} q_{i_k} F(t\ta_i) \right] \tag{\small{Convergence of $q_{i_k}$}}                                                         \\
				                                                    & = \lim_{k \rightarrow \infty} F^{-1} \left[  \sum_{i=1}^{N}  q_{i_k} F(t \ta_i) \right] \tag{$F^{-1}$ continuous \footnote{See \textcite[p. 337, Exercise 5.9.16]{elementaryAna}}} \\
				                                                    & =  \lim_{k \rightarrow \infty} t F^{-1} \left[  \sum_{i=1}^{N}  q_{i_k} F(\ta_i) \right] \tag{\cref{proof-homogen-rats}}                                                            \\
				                                                    & =  t F^{-1} \left[  \sum_{i=1}^{N}  w_i F(\ta_i) \right]. \tag{reversing the steps}
			\end{align*}
			Hence, the homogeneity holds for all $x_i, w_i, t \in \mathbb{R}_{\geq 0}$ and $N \in \mathbb{N}$, in particular, it holds for weights with sum $1$.
			By \cref{thm-power-mean-uniqueness} the only monotonic solutions are the power functions $\alpha x ^ p + \beta$ or the logarithm $\alpha \log(x)$ with $\alpha, \beta, p \in \mathbb{R}$ and $p > 0$. However, since $F(x) \geq 0$, it cannot be the logarithm and by $F(0) = 0$, the vertical shift $\beta$ must be equal to $0$. Therefore, the only solution is, that the extended function takes the form $\alpha x^p$, since we applied this theorem on the extension of $F$.
			Therefore, $F(x) = \alpha x^p$ must hold at least on the interval $[0, \Omega)$.
		\end{proofpp}
	\end{proofpart}


	\begin{proofpart}[Proving equality for positive reals]
		Now, for any $|\varphi_i(\ps{a}_i) - \varphi_i(\ps{b}_i)|$, we can find $t > 0$ small enough such that $t|\varphi_i(\ps{a}_i) - \varphi_i(\ps{b}_i)| < \omega$ for $i = 1, \ldots, n$. Applying monotonicity yields (we can assume $t \leq 1$):
		\begin{align*}
			\delta(\ps{a}, \ps{b}) & = F^{-1} \left( \sum_{i=1}^{n} F \left( |\varphi_i(\ps{a}_i) - \varphi_i(\ps{b}_i)| \right) \right)                                                                                      \\
			                       & = \frac{1}{t} F^{-1} \left( \sum_{i=1}^{n} F \left( t |\varphi_i(\ps{a}_i) - \varphi_i(\ps{b}_i)| \right) \right)                                                                        \\
			                       & = \frac{1}{t} \left( \sum_{i=1}^{n} \left( t|\varphi_i(\ps{a}_i) - \varphi_i(\ps{b}_i)| \right)^p \right)^{1/p} \tag{ since $t |\varphi_i(\ps{a}_i) - \varphi_i(\ps{b}_i)| \leq \omega$} \\
			                       & = \left( \sum_{i=1}^{n} |\varphi_i(\ps{a}_i) - \varphi_i(\ps{b}_i)|^p \right)^{1/p}                                                                                                      \\
			                       & = \metr[p]{\varphi_i(\ps{a}_i)}{\varphi_i(\ps{b}_i)}
		\end{align*}
		Thus, $\delta$ is a Minkowski metric and, by \cref{thm-repr-add-diff-struct}, the $F_i$ are interval scales, allowing for affine transformations. However, since we set $F_i(0)=0$, the affinity vanishes, and we are left with $F_i(x)= \alpha \cdot x^p$. Finally, the restriction $1 \leq p < \infty$ follows from Minkowski's inequality \parencite[pp. 189-191 Theorem 9]{handbookMeans}.
	\end{proofpart}
\end{proof}

